\documentclass[twocolumn, secnumarabic, amssymb, nobibnotes, aps, prd, amsmath, nofootinbib]{revtex4-2}

\newcommand{\revtex}{REV\TeX\ }
\newcommand{\classoption}[1]{\texttt{#1}}
\newcommand{\macro}[1]{\texttt{\textbackslash#1}}
\newcommand{\m}[1]{\macro{#1}}
\newcommand{\env}[1]{\texttt{#1}}
\setlength{\textheight}{9.5in}

\renewcommand{\topfraction}{0.99}
\renewcommand{\bottomfraction}{0.99}
\renewcommand{\textfraction}{0}
\setcounter{topnumber}{10}
\setcounter{bottomnumber}{10}
\setcounter{totalnumber}{10}

\DeclareMathOperator{\Tr}{Tr}
\DeclareMathOperator{\z}{\mathbf{z}}
\newcommand{\id}{\mathbbm{1}}
\newcommand{\cmark}{\ding{51}}%
\newcommand{\xmark}{\ding{55}}%
\newcommand{\SU}{\textrm{SU}}
%%%%%%%%%%%%%%%%%%%%%%%%%%%%%%%%%%%%%%%%%%%%%%%%%%%%%%%%%%%%%%%%%%%%%%%%%%%%%%%%
\newcommand{\bubble}{{\mathrm{b}}}
\newcommand{\metastable}{{\mathrm{m}}}
\newcommand{\stable}{{\mathrm{s}}}
\newcommand{\muca}{{\mathrm{muca}}}
\newcommand{\canonical}{{\mathrm{ca}}}
\newcommand{\weighted}{w}

\usepackage{graphicx}
\usepackage{hyperref}
\usepackage{xprintlen}
\usepackage{xcolor, algorithm2e}
\usepackage{soul}
\usepackage{mathtools}
\usepackage{svg}
\usepackage{subcaption}
\usepackage{amsmath, amssymb,amsfonts,amssymb,amscd} % For better math
\usepackage{bbm}
\usepackage[mode=buildnew]{standalone}% requires -shell-escape
%%%%%%%%%%%%%%%%%%%%%%%%%%%%%%%%%%%%%%%%%%%%%%%%%%%%%%%%%%%%%%%%%%%%%%%%%%%%%%%%

\usepackage{listings}
\usepackage{xcolor}

\lstset{
    language=C++,
    basicstyle=\ttfamily\small,
    keywordstyle=\color{blue}\bfseries,
    stringstyle=\color{red},
    commentstyle=\color{green!60!black},
    numbers=left,
    numberstyle=\tiny\color{gray},
    stepnumber=1,
    breaklines=true,
    frame=single,
    captionpos=b,
    showstringspaces=false
}
\begin{document}

\title{Resource contention in point to point overlapping GPU communication with MPI,NVSHMEM and NCCL}
\author{Aaron Haarti}
\email[]{aaron.haarti@helsinki.fi}

\affiliation{Department of Physics and Helsinki Institute of Physics,\\ PL64, FI-00014 University of Helsinki, Finland}

\date{Feb 2026}

\begin{abstract}
Concepts: P2P, NCCL, NVSHMEM, GPU-Aware MPI, Overlapping computation and communication
\end{abstract}

\maketitle

\section{Introduction}\label{sect:introduction}

\subsection{Communication Library Execution Models}

To accurately model resource contention, it is necessary to first define how standard communication libraries interface with the hardware. The architectural paradigms of MPI, NCCL, and NVSHMEM dictate their reliance on the GPU's Streaming Multiprocessors (SMs) and High-Bandwidth Memory (HBM), which fundamentally drives the contention profile.

\subsubsection{GPU-Aware MPI}
Modern implementations of the Message Passing Interface (MPI), such as MVAPICH2-GDR or OpenMPI, utilize Unified Virtual Addressing (UVA) and GPUDirect RDMA \cite{potluri2013mvapich2, nvidia_gpudirect}. In this host-initiated paradigm, the CPU orchestrates operations like \texttt{MPI\_Send} or \texttt{MPI\_Recv}. The Network Interface Card (NIC) accesses the GPU's HBM directly via the PCIe switch, entirely bypassing the GPU's SMs. Because the SMs are not involved in data staging or transmission, this model inherently isolates compute resources from communication traffic \cite{chu2020exploiting}.

\subsubsection{NVIDIA NCCL}
The NVIDIA Collective Communication Library (NCCL) operates on a host-initiated, kernel-executed model \cite{nvidia_nccl_doc}. According to NCCL architecture documentation, communication is managed by mapping thread blocks to SMs, referred to as "Channels." Standard NCCL routines (e.g., \texttt{ncclAllReduce}) employ a copy-based pipeline: the SMs actively copy data from the user buffer into internal FIFO staging buffers in the HBM, perform any necessary mathematical reductions, and then transmit the chunks to the NIC \cite{meta2025ncclx, unat2024landscape}. This approach requires active SM cycles and effectively doubles the HBM bandwidth requirement for a given payload \cite{agrawal2024optimizing}.

\subsubsection{NVIDIA NVSHMEM}
NVSHMEM implements a Partitioned Global Address Space (PGAS) designed for device-initiated communication \cite{nvidia_nvshmem_doc}. Unlike MPI or NCCL, APIs such as \texttt{nvshmem\_put} or \texttt{nvshmem\_get} are invoked directly from within the CUDA compute kernel. The GPU threads directly calculate destination addresses and ring the network doorbell to trigger the NIC \cite{unat2024landscape}. While this eliminates CPU orchestration overhead and delivers ultra-low latency for fine-grained messages, it requires the SMs to actively spend cycles issuing network requests rather than performing arithmetic computation \cite{hwang2023ark}. The exact SM involvement strictly depends on the transport layer: over intra-node NVLink, SMs perform active load/store instructions (high SM involvement); over InfiniBand (IBGDA), SMs simply write Work Queue Elements (WQEs) to the NIC (lower SM footprint); and on newer architectures like Hopper, the Tensor Memory Accelerator (TMA) handles the transfer autonomously, dropping SM involvement to near-zero \cite{nvidia_hopper_tma}. Despite variable SM usage, stream-enqueued \texttt{put} operations actively read data from the local HBM to send it. NVLink IPC transfers, for instance, can reach $\sim$75--80\% of the theoretical peak bandwidth, heavily saturating the shared HBM read ports. Consequently, overlapping NVSHMEM with memory-bound compute kernels still results in significant hardware contention \cite{nvidia_nvshmem_doc}.

\section{Resource Contention Analysis}

This analysis evaluates the performance impact of overlapping computation and communication. Using multiplicative degradation factors models resource contention realistically, as contention typically degrades the rate of execution, stretching task durations rather than adding a flat penalty.


Efficiency $E$ measures the sequential execution time relative to the overlapped execution time. Let $T_S$ be the total sequential time, $T_K$ be the computation time, $T_C$ be the communication time, and $T_O$ be the total overlapping time:
\begin{align*}
    E &= \frac{T_K + T_C}{T_O} \\
    T_K + T_C &\equiv T_S \implies E = \frac{T_S}{T_O}
\end{align*}

The loss function due to resource contention, $L$, accounts for the ratio between the theoretical perfectly overlapped time and the actual time including boundary computation $T_B$. Assuming boundary computation is zero ($T_{B} = 0$):
\begin{align*}
    L &= \frac{\max(T_K, T_C)}{T_O}
\end{align*}


\textbf{Ideal:} In a perfect system without contention, the degradation factors are 1. The loss function evaluates to 1, and the total time is strictly the maximum of computation or communication.
\begin{align*}
    L &= 1 \implies E = \frac{T_K + T_C}{\max(T_K, T_C)}
\end{align*}

\textbf{Contention:} When contention exists, it acts as a penalty factor that increases the total time of $T_K$ and $T_C$ individually. Let $a_{CL} \geq 1$ be the computation degradation factor and $b_{CL} \geq 1$ be the communication degradation factor. The total overlapping time $T_O$ becomes the maximum of these stretched times:
\begin{align*}
    T_O &= \max(T_K \cdot a_{CL}, T_C \cdot b_{CL})
\end{align*}
The degradation factors $a_{CL}$ and $b_{CL}$ are influenced by resource contention between the concurrent communication and compute kernels, depending on the RDMA engine, bandwidth, arithmetic intensity, and message size.

\subsection{Hardware-Dependent Formulation of Degradation Factors}

The degradation factors $a_{CL}$ and $b_{CL}$ quantify the structural bottlenecking of the GPU's Streaming Multiprocessors (SM) and High-Bandwidth Memory (HBM). Let $C_{SM}$ and $C_{HBM}$ denote the total hardware capacities for compute and memory bandwidth, respectively.

\subsubsection{Resource Demands}
A compute kernel $K$ with a specific arithmetic intensity $A$ demands resources $(R_{SM}^K, R_{HBM}^K)$.

A communication task $C$ handling message size $M$ demands resources $(R_{SM}^C, R_{HBM}^C)$. The magnitudes of $R_{SM}^C$ (software overhead and staging) and $R_{HBM}^C$ (physical DMA memory access) depend strictly on the transport mechanism utilized (e.g., copy-based versus zero-copy RDMA).

\subsubsection{Computation Degradation ($a_{CL}$)}
The computation degradation factor $a_{CL}$ is driven by SM thread starvation and HBM bandwidth saturation. It can be modeled as the maximum resource congestion ratio:
\begin{align*}
    a_{CL} &= \max\left(1, \frac{R_{SM}^K + R_{SM}^C}{C_{SM}}, \frac{R_{HBM}^K + R_{HBM}^C}{C_{HBM}}\right)
\end{align*}

\textbf{Impact of RDMA:} When utilizing zero-copy RDMA (e.g., GPU-Aware MPI or NCCLX), the communication task bypasses the SMs, yielding $R_{SM}^C = 0$. If the compute kernel is compute-bound (high $A$, requiring high $R_{SM}^K$ but low $R_{HBM}^K$), bypassing the SM prevents contention, ensuring $a_{CL} \approx 1$. Conversely, if the communication uses a copy-based protocol (e.g., standard NCCL), $R_{SM}^C > 0$, forcing $a_{CL} > 1$.

\subsubsection{Communication Degradation ($b_{CL}$)}
The communication degradation factor $b_{CL}$ represents the stretching of the network transfer time due to HBM contention throttling the NIC's injection rate. 
\begin{align*}
    b_{CL} &= \max\left(1, \frac{R_{HBM}^K + R_{HBM}^C}{C_{HBM}}\right)
\end{align*}

\textbf{Impact of Arithmetic Intensity:} If the overlapped compute kernel is memory-bound (low arithmetic intensity $A$), it demands massive HBM bandwidth ($R_{HBM}^K \to C_{HBM}$). This throttles the NIC's ability to fetch data via DMA from the GPU, forcing $b_{CL} > 1$ regardless of the RDMA engine's theoretical capacity.

\subsection{Methodology: Simulating Resource Contention}

To empirically measure the resource degradation factors—specifically the Streaming Multiprocessor contention ($a_{CL}$) and High-Bandwidth Memory contention ($b_{CL}$)—we developed a highly tunable benchmark kernel. Rather than executing a fixed application, this kernel acts as a synthetic stress test where the arithmetic intensity ($A$) can be dynamically modulated. This allows for the precise isolation of SM and HBM bottlenecks during overlapping communication.

\subsubsection{Modulating Arithmetic Intensity ($A$)}
The core mechanism controlling the kernel's behavior is the \texttt{math\_per\_load} parameter, which dictates the ratio of mathematical operations executed by the SMs relative to the memory fetches executed by the memory controller.
\begin{itemize}
    \item \textbf{Memory-Bound Workloads:} A low \texttt{math\_per\_load} forces the kernel to fetch data from HBM on nearly every loop iteration. This maximizes the compute kernel's memory demand ($R_{HBM}^K$), stressing the memory bandwidth and allowing the measurement of $b_{CL}$ when overlapped with DMA traffic.
    \item \textbf{Compute-Bound Workloads:} A high \texttt{math\_per\_load} executes hundreds of arithmetic instructions for every memory fetch. This maximizes SM demand ($R_{SM}^K$) while generating minimal memory traffic, allowing the measurement of $a_{CL}$ when overlapped with copy-based communication protocols.
\end{itemize}

\subsubsection{Simulating SM Contention}
To ensure genuine SM utilization, the kernel forces the execution of a continuous chain of dependent Fused Multiply-Add (FMA) instructions. Because the current value of the accumulator depends strictly on the immediate previous calculation, the compiler cannot unroll or optimize the loop away. This guarantees that the SMs are physically occupied for the specified number of iterations, creating a heavy compute footprint that will physically slow down if a communication operation attempts to steal SM cycles.

\subsubsection{Simulating HBM Contention}
To generate realistic memory pressure, the benchmark avoids simple contiguous memory reads, which the GPU's L1/L2 caches would easily absorb. Instead, it periodically reads from neighboring lattice sites based on 4D grid strides. By jumping through memory using varying strides, the kernel forces cache misses and guarantees that data must be pulled directly from the physical HBM. If a concurrent RDMA transfer is also pulling from HBM, the memory controller will throttle both tasks, exposing the true hardware memory bottleneck.

\subsubsection{Compute Kernel Implementation}
Listing \ref{lst:compute_kernel} details the benchmark compute kernel. The code is structured to enforce dependent mathematical operations while conditionally triggering strided memory accesses based on the target arithmetic intensity.

Listing \ref{lst:compute_kernel} details the benchmark compute kernel. The code is structured to enforce dependent mathematical operations while conditionally triggering strided memory accesses based on the target arithmetic intensity.

\begin{lstlisting}[language=C++, float=*t, caption={Benchmark compute kernel simulating arithmetic intensity and HBM contention.}, label={lst:compute_kernel}, basicstyle=\ttfamily\small, breaklines=true]
template <typename T>
__global__ void kernel_compute(T *d_bulk, const size_t *grid_size, size_t volume, 
                               int iters, int math_per_load) {
    size_t idx = blockIdx.x * (size_t)blockDim.x + threadIdx.x;
    if (idx >= volume) return;

    // Load grid dimensions into registers
    size_t g0 = grid_size[0], g1 = grid_size[1], g2 = grid_size[2];
    
    // Initial load
    T res = d_bulk[idx];
    T factor = (T)1.000001;

    // Strides to defeat simple contiguous caching
    size_t s0 = 1, s1 = g0, s2 = g0 * g1, s3 = g0 * g1 * g2;
    int load_count = 0;

    // HEAVY COMPUTE & MEMORY LOOP
    for (int i = 0; i < iters; ++i) {
        // 1. Dependent FMA ensures continuous SM utilization
        res = (res * factor) + (T)0.0123;
        
        if (res > 1e10 || res < -1e10) res *= 1e-10;

        // 2. Modulated HBM fetches based on Arithmetic Intensity
        int m_load = (math_per_load > 0) ? math_per_load : 1;
        if (i % m_load == 0) {
            size_t neighbor_idx;
            int mode = load_count % 4;
            
            if (mode == 0) neighbor_idx = (idx + s0) % volume;
            else if (mode == 1) neighbor_idx = (idx + s1) % volume;
            else if (mode == 2) neighbor_idx = (idx + s2) % volume;
            else neighbor_idx = (idx + s3) % volume;

            res += d_bulk[neighbor_idx];
            load_count++;
        }
    }
    // Final dependent write
    d_bulk[idx] = res;
}
\end{lstlisting}

\bibliography{bibliography}
\end{document}

\typeout{get arXiv to do 4 passes: Label(s) may have changed. Rerun}
